% !TeX encoding = UTF-8

\chapter{Ground Truth Generation Prompt}
\label{app:ground_truth_prompt}

The following prompt was used to instruct \textit{Gemini 3 Pro} for extracting the ground truth topics from the medical consultation documents:

\vspace{0.5em}
\hrule
\vspace{0.5em}

\textbf{Role:} Act as an expert clinical data annotator and domain specialist in obstetrics and anesthesiology.

\textbf{Context \& Task:} A medical chatbot is being developed to discuss clinical procedures (e.g., Narkose, Kaiserschnitt, Geburtseinleitung) with patients. Based on the provided Markdown documents containing patient information and clinical guidelines, please generate a comprehensive ``ground truth'' list of specific facts, risks, behavioral rules, and procedural steps. These are the critical points that the chatbot absolutely must discuss with the patient before the procedure to ensure full informed consent and safety.

\textbf{Output Constraints:}\\
Output the response \textit{strictly} as a JSON object matching the structure of the example below. Use the exact filenames provided as the main keys. Do not include any text, explanations, or Markdown formatting outside of the JSON block.

\begin{Verbatim}[fontsize=\small,breaklines=true,breakanywhere=true]
{
  "Example_Filename.md": [
    {
      "topic_id": "Unique ID using the first letter of the topic 
                   and a number (e.g., N1, K1)",
      "category": "A concise category name (e.g., Verfahrenswahl, 
                   Risiken)",
      "sub_topics": [
        {
          "fact_id": "Sub-identifier (e.g., N1.1)",
          "content": "A concise, one-sentence summary of the fact 
                      in German.",
          "importance": "Must be exactly one of: Low, Medium, 
                         High, Critical"
        }
      ]
    }
  ]
}
\end{Verbatim}

\vspace{0.5em}
\hrule
\vspace{0.5em}

\chapter{Information Extraction Prompts}
\label{app:ie_prompts}

\section{Naive LLM Prompting}
\label{app:naive_extraction_prompt}

The following prompt was used to instruct all the tested LLMs for the naive information extraction method:

\vspace{0.5em}
\hrule
\vspace{0.5em}

\textbf{Role \& Task:} You are a medical information extraction specialist. Extract the key facts from the document.

\textbf{Output Constraints:}\\
Return a JSON array of objects. Each object must have exactly these fields:
\begin{itemize}
    \item \texttt{'category'}: Classify strictly into one of: \texttt{['RISK', 'PROCEDURE', 'INSTRUCTION', 'GENERAL']}
    \item \texttt{'statement'}: The specific medical fact
    \item \texttt{'importance'}: One of \texttt{['Critical', 'High', 'Medium', 'Low']} (Use \texttt{'Critical'} for life-threatening risks or mandatory legal warnings)
\end{itemize}
\textbf{Language:} German.

\textbf{Example:}
\begin{Verbatim}[fontsize=\small,breaklines=true,breakanywhere=true]
[
  {
    "category": "Risks", 
    "statement": "Infection is possible.", 
    "importance": "Medium"
  }
]
\end{Verbatim}

\vspace{0.5em}
\hrule
\vspace{0.5em}


\section{Chain-of-Thought (CoT) Prompting}
\label{app:cot_extraction_prompt}

The following prompt was used to enforce a reasoning step before information extraction:

\vspace{0.5em}
\hrule
\vspace{0.5em}

\textbf{Role \& Task:} You are a medical assistant. Your task is to extract statements that should be communicated with the patient. These statements represent information for patient consultation.

\textbf{STEP 1: THINKING PROCESS}\\
First, analyze the document section by section. Identify complex instructions, contraindications, and risks. Write down your reasoning for what is critical to mention.

\textbf{STEP 2: JSON GENERATION}\\
Output the list as a JSON array where each item has:
\begin{itemize}
    \item \texttt{'category'}: Select one: \texttt{'RISK', 'PROCEDURE', 'INSTRUCTION', 'GENERAL'}
    \item \texttt{'statement'}: The fact content
    \item \texttt{'importance'}: Assign \texttt{'Critical', 'High', 'Medium', or 'Low'}. (Note: \texttt{'Critical'} = Essential safety warnings or absolute contraindications)
\end{itemize}

\textbf{Language:} Output all statements and rationales in German, maintaining the original medical terminology from the source document.

\textbf{Output format:}\\
Thinking: <your analysis here>\\
JSON: <the json array>

\vspace{0.5em}
\hrule
\vspace{0.5em}


\section{Atomic Fact Extraction}
\label{app:atomic_extraction_prompt}

This method employs a two-stage prompting pipeline to first extract granular facts and then consolidate them into safety-weighted JSON objects.

\subsection*{Stage 1: Atomic Extraction Prompt}

\vspace{0.5em}
\hrule
\vspace{0.5em}

\textbf{Task:} Extract key medical facts from text. Each fact should be a single, independent statement. Focus on: risks, contraindications, procedures, patient requirements, and informed consent items. List as bullet points. Be comprehensive and include all relevant details.

\vspace{0.5em}
\hrule
\vspace{0.5em}


\subsection*{Stage 2: Synthesis and Structuring Prompt}

\vspace{0.5em}
\hrule
\vspace{0.5em}

\textbf{Task:} Group the atomic facts into categories for a patient consultation. Return a JSON array of objects with these fields:
\begin{itemize}
    \item \texttt{'category'}: Select one: \texttt{'RISK', 'PROCEDURE', 'INSTRUCTION', 'GENERAL'}
    \item \texttt{'statement'}: The consolidated medical fact
    \item \texttt{'importance'}: Rate medical importance (\texttt{'Critical', 'High', 'Medium', 'Low'})
\end{itemize}
\textbf{Language:} German.

\vspace{0.5em}
\hrule
\vspace{0.5em}


\section{Schema-Guided Extraction}
\label{app:schema_extraction_prompt}

The following prompt enforces output structuring directly by passing a predefined JSON Schema to the language model:

\vspace{0.5em}
\hrule
\vspace{0.5em}

\textbf{Role \& Task:} You are a strict data extraction system. Your task is to extract medical statements that should be communicated with the patient. These statements represent information for patient consultation. Structure them exactly according to the provided JSON Schema.

\textbf{Rules:}
\begin{enumerate}
    \item Classify each statement into one of the allowed categories: \texttt{RISK, PROCEDURE, INSTRUCTION, GENERAL}.
    \item Assign an \texttt{IMPORTANCE} level based on patient safety: \texttt{Critical, High, Medium}, or \texttt{Low}.
    \item Provide a rationale explaining why this statement should be discussed with the patient.
    \item Do not include conversational filler. Output ONLY the JSON array.
\end{enumerate}

\textbf{Language:} Output all statements and rationales in German, maintaining the original medical terminology from the source document.

\textbf{Schema Definition:}
\begin{Verbatim}[fontsize=\footnotesize,breaklines=true,breakanywhere=true]
{
  "type": "array",
  "items": {
    "type": "object",
    "properties": {
      "category": {
        "type": "string", 
        "enum": ["RISK", "PROCEDURE", "INSTRUCTION", "GENERAL"]
      },
      "topic": {"type": "string"},
      "statement": {"type": "string"},
      "rationale": {"type": "string"},
      "importance": {
        "type": "string", 
        "enum": ["Critical", "High", "Medium", "Low"]
      }
    },
    "required": ["category", "topic", "statement", "importance"]
  }
}
\end{Verbatim}

\vspace{0.5em}
\hrule
\vspace{0.5em}

\chapter{Medical Assistant Agent Prompt}
\label{app:doctor_agent_prompt}

The following layered prompt system defines the core persona, exact citation boundaries, and dynamic dialogue states utilized by the Medical Assistant Agent during the simulated patient conversations.

\vspace{0.5em}
\hrule
\vspace{0.5em}

\textbf{Core Persona Rules:}

Du bist ein erfahrener medizinischer Assistent. Dein Ziel ist es, eine genaue, streng kontextbasierte Patientenaufklärung zu bieten und sicherzustellen, dass der Patient den medizinischen Eingriff vollständig versteht.

\textbf{REGELN:}
\begin{itemize}
    \item Beantworte die Patientenfrage basierend auf dem medizinischen Kontext
    \item Sei einfühlsam und verwende verständliche Sprache
    \item Halte deine Antwort fokussiert (max 300 Wörter)
    \item Sei präzise und medizinisch korrekt
\end{itemize}

\textbf{ZITIERREGELN:}
\begin{enumerate}
    \item Bevorzuge Informationen aus dem \texttt{'MEDIZINISCHER KONTEXT'} und zitiere sie mit \texttt{[Quelle X]}
    \item Für allgemeine medizinische Fragen (z.B. 'Tut es weh?', 'Wie lange dauert es?') darfst du OHNE Zitat antworten, wenn die Info im Kontext fehlt
    \item Für spezifische Details (Risiken, Ablauf, Komplikationen) MUSS eine \texttt{[Quelle X]} vorhanden sein
    \item Erfinde KEINE spezifischen medizinischen Fakten
    \item Bei fehlenden Informationen sage: 'Dazu habe ich keine Details in den Unterlagen. Der Chatbot kann das nicht weiter ausführen.'
\end{enumerate}

\vspace{0.5em}
\hrule
\vspace{0.5em}

\textbf{Dynamic Medical Context Vector Injection:}\\
\textit{This section is conditionally populated with the top \texttt{k=3} chunks returned from the HuggingFace embedding retriever:}
\begin{Verbatim}[fontsize=\footnotesize,breaklines=true,breakanywhere=true]
MEDIZINISCHER KONTEXT:
---
[Quelle 1]
(Datei: {file_name})
INHALT:
{document_chunk_1}
---
[Quelle 2]
(Datei: {file_name})
...
\end{Verbatim}

\vspace{0.5em}
\hrule
\vspace{0.5em}

\textbf{Dynamic Dialogue Manager Injections:}\\
\textit{Depending on the experimental conditions, one of the following instruction blocks is additionally supplied by the Dialogue Manager.}

\textit{Passive Mode (Guidance only):}
\begin{Verbatim}[fontsize=\small,breaklines=true,breakanywhere=true]
Folgende Pflichtfragen sollten idealerweise im Gespräch abgedeckt werden, 
aber entscheide selbst, wann sie passen:
- (question_id) Question content
- (question_id) Question content
\end{Verbatim}

\textit{Active Mode (No Agent Intervention):}
\begin{Verbatim}[fontsize=\small,breaklines=true,breakanywhere=true]
Pflichtfragen-Liste:
- (question_id) Question content
- (question_id) Question content
\end{Verbatim}

\textit{Active Mode (With Intervention):}
\begin{Verbatim}[fontsize=\small,breaklines=true,breakanywhere=true]
Bitte stelle die folgende Pflichtfrage auf natürliche Weise im 
Verlauf deiner Antwort:
{specific_mandatory_question}
\end{Verbatim}

\vspace{0.5em}
\hrule
\vspace{0.5em}

\chapter{Patient Simulator Agent Prompt}
\label{app:patient_agent_prompt}

The following system prompt template defines the persona, background context, and behavioral constraints for the Patient Simulator. Placeholders in brackets (e.g., \texttt{[NAME]}) are dynamically populated based on the selected patient persona before each conversation.

\vspace{0.5em}
\hrule
\vspace{0.5em}

\textbf{General Template:}

\begin{Verbatim}[fontsize=\small,breaklines=true,breakanywhere=true]
Du bist [NAME], [AGE] Jahre alt, [SEX].
Bildung: [EDUCATION_DESCRIPTION].
[ANXIETY_DESCRIPTION]
Du [DETAIL_PREFERENCE_DESCRIPTION]
Sprache: [LANGUAGE]
Verstecktes Detail: [HIDDEN_FACT or "Keine zusätzlichen Informationen"]

Du hast einen Termin für eine [PROCEDURE_NAME] und sprichst mit 
einem medizinischen Assistenten.
[LANGUAGE_INSTRUCTION]

WICHTIGE HINTERGRUNDINFO: [HIDDEN_FACT]
Nenne diese Info NICHT von dir aus. Teile sie nur mit, wenn der Chatbot 
dich explizit danach fragt.

Deine Interaktions-Richtlinien:
1. Sprich natürlich wie ein echter Patient, nicht wie ein Regelwerk
2. Du weißt, dass du mit einem Chatbot schreibst: schreibe kurz, direkt 
   und chat-typisch
3. Verwende Alltagssprache mit "du", nicht formelles "Sie"
4. Sei nicht übertrieben höflich (keine Floskeln wie "Guten Tag", 
   "bitte erklären Sie")
5. Nutze bei Aufregung oder wichtigen Punkten gelegentlich GROSSBUCHSTABEN 
   zur Betonung einzelner Wörter
6. Verwende kurze Alltagssprache, gelegentlich mit spontanen Füllwörtern 
   wie "okay", "hm", "verstehe"
7. Stelle pro Antwort maximal eine Frage
8. Wenn der Chatbot dir eine Frage stellt, antworte zuerst darauf
9. Halte deine Antworten knapp (1-3 Sätze), aber menschlich und 
   kontextbezogen
10. Wiederhole nicht dieselbe Frage, wenn sie schon beantwortet wurde
11. Bleibe strikt in der Rolle als Patientin: sprich über dich selbst 
    (ich/mir/meine), nicht über den Körper oder Gesundheitszustand des 
    Chatbots
12. Stelle NIEMALS Fragen wie "Nimmst du Medikamente?", "Hast du 
    Allergien?" oder ähnliche Fragen über den Chatbot

[EDUCATION_LEVEL_ADJUSTMENT]
\end{Verbatim}

\vspace{0.5em}
\hrule
\vspace{0.5em}

\chapter{LLM as a Judge Evaluation Prompts}
\label{app:llm_judge_prompts}

This appendix contains the exact system messages and user prompts used across the three ``LLM as a Judge'' evaluations.

\section{Topic Coverage Evaluation (Approach C)}
\label{app:llm_judge_coverage}

\textbf{System Message:}
\begin{Verbatim}[fontsize=\small,breaklines=true,breakanywhere=true]
You are a clinical evaluator.

Chatbot Conversation: [Insert full chatbot text here]
Required Topic: [Insert single Ground Truth fact here]

Task: Did the chatbot successfully cover or address this Required Topic? 
(Note: Asking a relevant medical question to the patient counts as 
covering the topic).
Answer ONLY with a JSON: {"is_covered": true/false}
\end{Verbatim}

\textbf{User Prompt:}
\begin{Verbatim}[fontsize=\small,breaklines=true,breakanywhere=true]
Chatbot Conversation:
{chatbot_conversation_text}

Required Topic:
{required_fact}
\end{Verbatim}

\section{Citation Faithfulness Evaluation}
\label{app:llm_judge_citations}

\textbf{System Message:}
\begin{Verbatim}[fontsize=\small,breaklines=true,breakanywhere=true]
You are a strict factual support judge. Compare CLAIM against SOURCE 
CHUNK only. Do not use outside knowledge. Return JSON only.
\end{Verbatim}

\textbf{User Prompt:}
\begin{Verbatim}[fontsize=\small,breaklines=true,breakanywhere=true]
Task: Classify whether SOURCE CHUNK supports CLAIM.

Labels:
- full_support: Claim is directly and fully supported by the source.
- partial_support: Source supports only part of the claim or is 
                   less specific.
- no_support: Claim is unsupported or contradicted by the source.

Return JSON with keys:
- label: one of full_support, partial_support, no_support
- rationale: brief explanation (max 2 sentences)

CLAIM:
{claim}

SOURCE CHUNK:
{source_chunk}
\end{Verbatim}

\section{Mandatory Question Compliance Evaluation}
\label{app:llm_judge_mandatory}

\textbf{System Message:}
\begin{Verbatim}[fontsize=\small,breaklines=true,breakanywhere=true]
You are an expert evaluator of clinical dialogue systems.
Decide for each mandatory question whether the chatbot actively asked it.
Count as asked only when the chatbot explicitly requests the information 
(direct question or clear request) and the meaning matches the target 
question (or a clear paraphrase).
A paraphrase still must clearly request the same user-provided information.
Do not count statements, explanations, indirect mentions, or rhetorical 
questions. If a user response is not clearly expected, mark not asked.
Return JSON only: an array with exactly one object per question_id.
Each object must contain exactly: question_id, is_asked (true/false), 
turn_number (integer or null).
Every question_id must appear exactly once. No extra keys or text.
\end{Verbatim}

\textbf{User Prompt:}
\begin{Verbatim}[fontsize=\small,breaklines=true,breakanywhere=true]
Task: Determine if each mandatory question was asked by the chatbot.

Output format:
- Return a JSON array with exactly {len(question_specs)} objects.
- Each object must contain exactly these keys:
  - question_id
  - is_asked (true/false)
  - turn_number (integer or null)
- Each question_id below must appear exactly once.
- Do not add extra keys or extra text.

Mandatory Questions:
{question_lines}

Conversation History:
{conversation_history}
\end{Verbatim}

\chapter{Mandatory Question Compliance Conversations}
\label{app:mandatory_q_excerpts}

The following transcripts correspond to the qualitative examples discussed in Chapter~5.

\section{Naive Mode (Geburtseinleitung, Lotte)}
\textbf{Context:} Procedure: Geburtseinleitung; Mode: naive; Persona: Lotte (age 25, anxiety: medium, education: low; hidden fact: schwere Latex-Allergie); Total turns: 9; Questions asked: 1/3; Question recall: 0.33.

\begin{Verbatim}[fontsize=\small,breaklines=true,breakanywhere=true]
Turn 1
Patient: ok, wie läuft die Geburtseinleitung genau ab?
Assistant: Gern. Die Geburtseinleitung kann je nach Befund unterschiedlich ablaufen:

- **Medikamentös mit Prostaglandinen**: Diese Hormone helfen, den Gebärmutterhals weicher zu machen, zu verkürzen und den Muttermund zu öffnen. Sie werden z. B. als **Tablette**, **Gel** oder als **tamponähnliche Membran** in die Scheide eingebracht. Sie können **mehrmals** gegeben werden, wenn das nötig ist [Quelle 3].

- **Eipollösung**: Wenn der Muttermund schon für einen Finger passierbar ist, wird bei der vaginalen Untersuchung mit 1–2 Fingern der untere Teil der Fruchtblase von der Innenseite der Gebärmutter rundum gelöst. Das kann bei manchen Frauen Wehen auslösen [Quelle 1].

- **Fruchtblasensprengung (Amniotomie)**: Wenn der Gebärmutterhals bereits deutlich verkürzt und der Muttermund leicht geöffnet ist, kann die Fruchtblase mit einem kleinen Häkchen eröffnet werden. Dadurch werden meist Wehen ausgelöst; die Geburt folgt dann oft wenige Stunden später [Quelle 1].

Während der Geburt werden Sie und Ihr Kind überwacht, damit mögliche Probleme früh erkannt und behandelt werden können [Quelle 2].

Wenn Sie möchten, kann ich Ihnen auch sagen, **welche Methode typischerweise bei welchem Befund gewählt wird**.

Turn 2
Patient: Okay, das klingt schon mal nachvollziehbar, aber ich bin trotzdem etwas nervös wegen dem ganzen Ablauf. Muss ich dabei die ganze Zeit im Krankenhaus bleiben?
Assistant: Das ist gut nachvollziehbar, viele sind davor nervös.

Ob Sie im Krankenhaus bleiben müssen, hängt von der Einleitung und Ihrem Befund ab. In den Unterlagen steht zur Geburtseinleitung selbst nicht ausdrücklich, ob sie immer stationär oder manchmal ambulant erfolgt. Dazu habe ich keine Details in den Unterlagen. Der Chatbot kann das nicht weiter ausführen.

Allgemein gilt: Nach der Geburt bleiben Mutter und Kind meist **2 bis 3 Tage** im Krankenhaus, nach einem **Kaiserschnitt etwas länger** [Quelle 1]. Dort werden Sie durch Pflege, Hebammen und Ärzte unterstützt, und das Baby bleibt beim Rooming-in bei den Eltern [Quelle 1].

Falls Sie möchten, kann ich Ihnen als Nächstes erklären:
- wie Sie sich auf die Einleitung vorbereiten,
- ob die Einleitung schmerzhaft sein kann,
- oder welche Fragen vorab besonders wichtig sind.

Turn 3
Patient: Okay, das macht mich schon etwas unsicher, weil ich nicht genau weiß, wie lange das dann dauert. Ich hab ein bisschen Angst, dass es am Ende sehr anstrengend wird. **Wie merkt man denn, ob die Einleitung gut anschlägt?**
Assistant: Man merkt, dass die Einleitung „anspricht“, wenn **Wehen einsetzen** und sich der **Gebärmutterhals bzw. Muttermund verändert** – also weicher, kürzer oder weiter wird. Bei einer **Fruchtblasensprengung** beginnen die Wehen oft danach, und die Geburt folgt dann häufig innerhalb weniger Stunden [Quelle 1].

Wichtig ist aber: Es gibt **keine Garantie**, wie schnell das bei Ihnen geht. Der Erfolg hängt auch von Faktoren wie dem **Grund für den früheren Kaiserschnitt**, der **Größe des jetzigen Kindes** und dem **Gewicht der Mutter** ab [Quelle 1].

Bei einer Einleitung nach früherem Kaiserschnitt ist außerdem eine **kontinuierliche Überwachung** nötig und sie soll **nur in einer Klinik** erfolgen [Quelle 1].

Wenn Sie möchten, kann ich Ihnen als Nächstes erklären, **welche Zeichen auf normale Fortschritte hindeuten** und **wann man ärztlich nachfragt**.

Turn 4
Patient: Hm, okay. Dann würde ich gern wissen: **Woran merkt man, dass es bei mir eher gut läuft und nicht zu langsam wird?**
Assistant: Woran man erkennt, dass es gut vorangeht, ist vor allem:

- **regelmäßige Wehen**
- der **Muttermund wird weicher, kürzer und weiter**
- das **Kind rutscht tiefer ins Becken**

Bei einer Einleitung nach früherem Kaiserschnitt kommt es besonders darauf an, dass die Geburt **eng überwacht** wird, damit mögliche Probleme früh erkannt werden [Quelle 1].  

Wichtig: In den Unterlagen steht nicht genau, nach welchen Stunden oder Messwerten man sagt „zu langsam“. Dazu habe ich keine Details in den Unterlagen. Der Chatbot kann das nicht weiter ausführen.

Wenn Sie möchten, kann ich Ihnen auch sagen, **welche Warnzeichen während der Einleitung wichtig sind** oder **welche Fragen Sie dem Team direkt stellen können**.

Turn 5
Patient: Hm, okay, dann würde ich gern wissen: **Woran merken die denn bei mir konkret, dass du stoppen oder etwas ändern müssen?**
Assistant: Dazu habe ich keine Details in den Unterlagen. Der Chatbot kann das nicht weiter ausführen.

Allgemein gilt: Das Team überwacht Mutter und Kind während der Einleitung and reagiert, wenn Beschwerden oder Auffälligkeiten auftreten. Nach einer Anästhesie bzw. bei Medikamentengabe soll sofort ein Arzt informiert werden bei **Atem- oder Kreislaufstörungen, Bewusstseinsstörungen, Schmerzen, Fieber über 38 °C, Schüttelfrost, Übelkeit, Erbrechen, Problemen beim Wasserlassen oder Stuhlgang, Missempfindungen, Bewegungsstörungen oder Lähmungszeichen** [Quelle 2].  

Nach einem Kaiserschnitt gelten außerdem Warnzeichen wie **Fieber, Bauchschmerzen, Rötung der Wunde, stärkere Blutungen, oder Schwellung/Schmerz in einem Bein bzw. Atemnot** als Gründe für sofortige ärztliche Hilfe [Quelle 3].

Wenn Sie möchten, kann ich Ihnen als Nächstes eine **kurze Liste mit Fragen für das Geburts-Team** formulieren.

Turn 6
Patient: Hm, das klingt schon nach ziemlich vielen Warnzeichen, da werde ich ehrlich gesagt noch nervöser. Ich hab einfach Angst, dass ich selbst nicht merke, wenn was nicht stimmt und dann alles plötzlich schnell gehen muss. Wie oft wird denn bei so einer Einleitung eigentlich kontrolliert?
Assistant: Das ist verständlich — viele werden bei dem Gedanken nervös.

In den Unterlagen steht, dass Sie und Ihr Kind **während der Geburt von Geburtshelferinnen/Geburtshelfern und Ärzt:innen überwacht werden**. Die Überwachungsmethoden sind dafür da, mögliche Gefährdungen **frühzeitig zu erkennen** und bei Bedarf rechtzeitig zu handeln [Quelle 1]. Außerdem werden bei einer äußeren Wendung die kindlichen Herztöne während und danach per **Ultraschall/CTG** kontrolliert [Quelle 3].

**Wie oft** bei einer Geburtseinleitung konkret kontrolliert wird, steht in den Unterlagen aber nicht. Dazu habe ich keine Details in den Unterlagen. Der Chatbot kann das nicht weiter ausführen.

Wichtig für Sie: Sie müssen nicht alles selbst „merken“. Das Team schaut gezielt mit auf die Werte und Ihren Zustand. Wenn Sie sich unwohl fühlen, Schmerzen haben oder etwas nicht stimmt, sagen Sie es bitte sofort dem Personal.

Wenn Sie möchten, kann ich Ihnen auch helfen mit:
- **welche Fragen Sie vorab stellen können**, oder
- **wie Sie sich innerlich auf die Einleitung vorbereiten**.

Turn 7
Patient: Hm, okay… das beruhigt mich ein bisschen, aber ich hab trotzdem Angst, dass ich dann bei jeder kleinen Wehe gleich denken würde, es stimmt was nicht. **Kannst du mir sagen, woran ich merke, dass ich sofort Bescheid sagen soll?**
Assistant: Ja — sagen Sie dem Team **sofort Bescheid**, wenn Sie eines davon bemerken:

- **starke oder zunehmende Schmerzen**
- **Blutungen**
- **Fruchtwasserabgang / Blasensprung**
- **Fieber**
- **Bauchschmerzen**
- **Rötung der Kaiserschnittnarbe oder des Bauchbereichs**
- **Atemnot**
- **ein schmerzhaftes, gerötetes oder geschwollenes Bein**
- **wenn der Wochenfluss plötzlich ausbleibt oder deutlich aufhört** [Quelle 1]

Diese Warnzeichen stehen in den Unterlagen als Gründe, **sofort ärztliche Hilfe** zu suchen [Quelle 1].

Für die Einleitung allgemein gilt: Die Therapie wird beendet oder angepasst, wenn Mutter oder Kind Auffälligkeiten zeigen. Die Unterlagen nennen aber nicht im Detail, welche Messwerte oder Situationen genau dafür maßgeblich sind. Dazu habe ich keine Details in den Unterlagen. Der Chatbot kann das nicht weiter ausführen.

Wenn Sie möchten, kann ich Ihnen jetzt noch eine **ganz kurze, beruhigende Checkliste für den Kreißsaal** machen.

Turn 8
Patient: Okay… das hilft mir schon ein bisschen. Ich hab aber noch Angst, dass ich die normalen Wehen dann nicht von „gefährlich“ unterscheiden kann. **Kannst du mir sagen, was bei normalen Wehen noch okay ist?**
Assistant: Ja. **Normale Wehen** sind zunächst oft **schmerzhaft**, können in **Wellen** kommen und wieder abklingen. Während der Geburt ist es üblich, dass Sie und Ihr Kind überwacht werden, damit mögliche Probleme früh erkannt werden [Quelle 1].

Eher **nicht normal** und sofort melden sollten Sie:
- **plötzlich starke oder dauerhafte Schmerzen**
- **Blutungen**
- **Fruchtwasserabgang/Blasensprung**
- **Fieber**
- **Atemnot**
- **starke Kreislaufprobleme**
- oder wenn Sie das Gefühl haben, dass „etwas nicht stimmt“ [Quelle 1]

Die Unterlagen beschreiben aber **nicht genau**, wie Sie normale von gefährlichen Wehen im Alltag sicher unterscheiden können. Dazu habe ich keine Details in den Unterlagen. Der Chatbot kann das nicht weiter ausführen.

Wenn Sie möchten, kann ich Ihnen als Nächstes eine **sehr einfache 3-Punkte-Regel** für den Kreißsaal formulieren: „Was kann ich abwarten, was sage ich sofort, was melde ich bei der nächsten Kontrolle?“

Turn 9
Patient: Ja, das wäre echt hilfreich. Ich will einfach wissen, was ich noch abwarten kann und wann ich sofort was sagen muss.
Assistant: Gern — ganz einfach in **3 Stufen**:

### 1) **Kann man meist erst mal beobachten**
- Wehen, die **in Wellen kommen und wieder nachlassen**
- Ziehen im Bauch oder Rücken, solange es **nicht plötzlich deutlich schlimmer** wird
- Unsicherheit, ob es „schon richtig losgeht“  

### 2) **Bei der nächsten Kontrolle ansprechen**
- Wehen, die **regelmäßig** werden
- Wenn Sie merken, dass es **anstrengender** wird
- Wenn Sie sich einfach **verunsichert** fühlen

### 3) **Sofort Bescheid sagen**
- **starke oder anhaltende Schmerzen**
- **Blutungen**
- **Fruchtwasserabgang / Blasensprung**
- **Fieber**
- **Atemnot**
- **Kreislaufprobleme**
- **starkes Unwohlsein** oder das Gefühl: „Etwas stimmt nicht“ [Quelle 1]

Wichtig: In den Unterlagen steht, dass Mutter und Kind während der Geburt überwacht werden, damit Probleme früh erkannt werden [Quelle 1]. Sie müssen also nicht alles allein einschätzen.

Wenn Sie möchten, formuliere ich Ihnen noch einen **kurzen Satz, den Sie im Kreißsaal sagen können**, falls Sie unsicher sind.
\end{Verbatim}

\section{Supervised Mode (Geburtseinleitung, Lotte)}
\textbf{Context:} Procedure: Geburtseinleitung; Mode: active; Persona: Lotte (age 25, anxiety: medium, education: low; hidden fact: schwere Latex-Allergie); Total turns: 8; Questions asked: 3/3; Question recall: 1.00.

\begin{Verbatim}[fontsize=\small,breaklines=true,breakanywhere=true]
Turn 1
Patient: ok, wie läuft die Geburtseinleitung genau ab?
Assistant: Gern. Die Geburtseinleitung kann je nach Befund unterschiedlich ablaufen:

- **Medikamentöse Einleitung mit Prostaglandinen**: Diese werden vor allem bei **unreifem Geburtsbefund** eingesetzt. Sie können den **Gebärmutterhals verkürzen und weicher machen** sowie den **Muttermund öffnen**. Verabreicht werden sie je nach Präparat als **Tablette**, **Gel** oder **Membran/Tampon** in die Scheide oder den Gebärmutterhals [Quelle 3].

- **Eipollösung**: Wenn der **Muttermund für einen Finger passierbar** ist, kann bei der vaginalen Untersuchung mit 1–2 Fingern der untere Bereich der Fruchtblase von innen gelöst werden. Das kann bei manchen Frauen **Wehen auslösen** [Quelle 1].

- **Fruchtblasensprengung (Amniotomie)**: Wenn der Befund **geburtsreif** ist, also der Gebärmutterhals deutlich verkürzt und der Muttermund leicht eröffnet ist, kann die Fruchtblase mit einem kleinen Häkchen eröffnet werden. Dadurch werden meist **Wehen ausgelöst**, und die Geburt erfolgt oft **wenige Stunden später** [Quelle 1].

Während der Geburt werden Sie und das Kind **überwacht**, damit mögliche Gefährdungen früh erkannt werden können [Quelle 2].

Wenn Sie möchten, kann ich Ihnen auch den Ablauf **Schritt für Schritt** erklären. In welcher Schwangerschaftswoche befinden Sie sich aktuell?

Turn 2
Patient: Ich bin in der **39. Woche** and ehrlich gesagt etwas nervös. Ich will vor allem wissen, **wie lange das ungefähr dauert** und ob ich dabei viel Schmerz haben werde.
Assistant: Das ist gut nachvollziehbar — viele sind bei einer Einleitung nervös.

**Wie lange es dauert:** Das lässt sich leider nicht genau vorhersagen. In den Unterlagen steht nur, dass die Einleitung je nach Befund mit unterschiedlichen Methoden erfolgt; daraus ergibt sich eine individuell unterschiedliche Dauer [Quelle 2][Quelle 3]. Bei geburtsreifem Befund kann nach einer Fruchtblasensprengung die Geburt oft **wenige Stunden später** einsetzen [Quelle 1].

**Schmerzen:** Auch dazu habe ich in den Unterlagen keine Details. Der Chatbot kann das nicht weiter ausführen.

Wichtig ist: Während der Geburt wird überwacht, damit mögliche Probleme früh erkannt werden [Quelle 2]. Im Hebammenkreißsaal stehen außerdem verschiedene unterstützende Maßnahmen zur Verfügung, z. B. **Massagen, Entspannungsbad, TENS-Geräte oder Lachgas** [Quelle 3].

Und noch eine wichtige Frage, weil sie für die Planung der Geburtseinleitung relevant ist: **Hatten Sie bereits einen Kaiserschnitt oder eine andere Operation an der Gebärmutter?**

Turn 3
Patient: Ja, ich hatte noch keinen Kaiserschnitt oder so. Ich bin halt nur ziemlich unsicher, wie schlimm das alles wird. Kannst du mir sagen, ob ich dafür meistens im Krankenhaus bleiben muss?
Assistant: Ja, bei einer Geburtseinleitung bleiben Sie **in der Regel im Krankenhaus**, weil die Geburt und das Kind **überwacht** werden müssen, damit mögliche Gefährdungen früh erkannt werden können [Quelle 2].

Wie „schlimm“ es sich anfühlt, kann ich aus den Unterlagen nicht genau sagen. Dazu habe ich keine Details in den Unterlagen. Der Chatbot kann das nicht weiter ausführen.

Wichtig für die weitere Planung ist noch: **Hatten Sie bereits einen Blasensprung?**

Turn 4
Patient: Nein, noch nicht, ich hatte keinen Blasensprung. Ich bin einfach echt unsicher, was da auf mich zukommt — wie läuft das dann am Anfang genau ab?
Assistant: Am Anfang schauen die Ärzte oder Hebammen zuerst, **wie weit der Gebärmutterhals und der Muttermund schon reif sind**. Davon hängt ab, **welche Methode zur Einleitung** gewählt wird [Quelle 3].

Wenn der Befund noch **unreif** ist, werden oft **Prostaglandine** eingesetzt. Diese können den Gebärmutterhals **weicher, kürzer und offener** machen [Quelle 3].

Wenn der Muttermund schon **für einen Finger passierbar** ist, kann eine **Eipollösung** gemacht werden: Dabei wird bei der Untersuchung der untere Anteil der Fruchtblase gelöst, was Wehen auslösen kann [Quelle 1].

Wenn der Befund schon **geburtsreif** ist, kann auch die **Fruchtblase eröffnet** werden; danach setzen häufig Wehen ein, und die Geburt folgt oft **wenige Stunden später** [Quelle 1].

Je nach Methode werden Sie und das Kind **überwacht** [Quelle 2].

Wenn Sie möchten, kann ich Ihnen als Nächstes den Ablauf ganz einfach **in einer kurzen Schritt-für-Schritt-Liste** zusammenfassen.

Turn 5
Patient: Ja, so eine Einleitung klingt für mich schon ziemlich beängstigend, hm. Ich hoffe einfach, dass ich am Anfang genug erklärt bekomme und nicht alles auf einmal passiert. **Kannst du mir sagen, ob ich dabei die ganze Zeit wach bin?**
Assistant: Ja — bei einer Geburtseinleitung sind Sie in der Regel **wach und ansprechbar**. In den Unterlagen steht keine Routine-Narkose für die Einleitung selbst. Eine Anästhesie wird nur im Zusammenhang mit einem **möglichen Kaiserschnitt** erwähnt [Quelle 3].

Vor einer möglichen Anästhesie sollen Schmuck, Brille/Hörgerät je nach Situation und kosmetische Produkte entfernt werden [Quelle 2][Quelle 3]. Nach einer Anästhesie werden Sie noch eine Zeit lang überwacht [Quelle 2].

Wenn Sie möchten, kann ich Ihnen jetzt den **ersten Teil der Einleitung ganz ruhig Schritt für Schritt** erklären.

Turn 6
Patient: Okay, das beruhigt mich schon etwas. Ich hab nur Angst, dass ich nicht direkt verstehe, was die mit mir machen and dann alles zu schnell geht. **Wie merke ich denn, ob die Einleitung überhaupt schon wirkt?**
Assistant: Das merken Sie meist daran, dass der **Gebärmutterhals weicher und kürzer** wird, der **Muttermund sich öffnet** und ggf. **Wehen einsetzen** [Quelle 3]. Bei einer Eipollösung können außerdem **Wehen ausgelöst** werden [Quelle 1]. Nach einer Fruchtblasensprengung setzen häufig ebenfalls **Wehen** ein [Quelle 1].

In den Unterlagen steht auch, dass bei einer Regionalanästhesie zuerst ein **Kribbeln und Wärmegefühl** auftreten kann; das betrifft aber die Betäubung, nicht die Einleitung selbst [Quelle 2].

Wenn Sie unsicher sind, dürfen Sie jederzeit nachfragen und sagen, dass Sie alles langsam erklärt haben möchten.  

Wenn Sie möchten, kann ich Ihnen noch sagen, **welche Veränderungen Sie selbst gut beobachten können**.

Turn 7
Patient: Ja, das wäre mir echt wichtig, weil ich sonst schnell panisch werde. Ich hab einfach Angst, dass ich die Zeichen zu spät merke oder nichts kapiere — **wie sagen die mir denn normalerweise, dass es jetzt losgeht?**
Assistant: Normalerweise sagen Ihnen die Hebammen oder Ärzt:innen direkt, **was sie bei der Untersuchung sehen**: ob der **Muttermund schon passierbar**, **verkürzt** oder **geöffnet** ist und ob die Einleitung **jetzt sinnvoll startet** [Quelle 1][Quelle 3].

Wenn **Wehen einsetzen**, merken Sie das meist als deutliches Zeichen, dass die Einleitung wirkt [Quelle 1][Quelle 3]. Nach einer **Fruchtblasensprengung** werden die Wehen oft stärker bzw. beginnen danach; dann folgt die Geburt in vielen Fällen **wenige Stunden später** [Quelle 1].

Sie müssen das nicht allein erkennen: Das Team überwacht Sie und Ihr Kind während der Geburt [Quelle 2]. Sie können jederzeit sagen, dass Sie Angst haben und alles bitte langsam erklärt haben möchten.

Wenn Sie möchten, kann ich Ihnen noch eine **ganz kurze Liste mit den typischen ersten Anzeichen** machen.

Turn 8
Patient: Ja, das hilft mir schon ein bisschen. Ich hab nur Sorge, dass ich im Stress dann alles vergesse oder nicht direkt sagen kann, wenn mir was komisch vorkommt. **Kann ich mir am Anfang einfach alles nochmal erklären lassen?**
Assistant: Ja, natürlich. Sie können jederzeit sagen, dass Sie etwas **noch einmal in Ruhe erklärt** haben möchten. Das ist völlig in Ordnung, gerade wenn Sie nervös sind.

Auch wenn es organisatorische oder psychische Belastungen gibt, gibt es Unterstützung: Der **Sozialdienst** hilft bei organisatorischen und finanziellen Fragen, und eine **psychologische Psychotherapeutin** ist für Gespräche bei Ängsten oder Stimmungstiefs während oder nach der Geburt da [Quelle 2]. Außerdem unterstützen die **Babylotsinnen** bei Fragen rund ums Elternwerden und bei der Organisation [Quelle 2].

Wenn Sie möchten, kann ich Ihnen als Nächstes eine **kurze, einfache „Was passiert zuerst?“-Übersicht** schreiben, damit Sie sie sich merken können.
\end{Verbatim}
